\IfFileExists{stacks-project.cls}{%
\documentclass{stacks-project}
}{%
\documentclass{amsart}
}

% For dealing with references we use the comment environment
\usepackage{verbatim}
\newenvironment{reference}{\comment}{\endcomment}
%\newenvironment{reference}{}{}
\newenvironment{slogan}{\comment}{\endcomment}
\newenvironment{history}{\comment}{\endcomment}

% For commutative diagrams we use Xy-pic
\usepackage[all]{xy}

% We use 2cell for 2-commutative diagrams.
\xyoption{2cell}
\UseAllTwocells

% We use multicol for the list of chapters between chapters
\usepackage{multicol}

% This is generally recommended for better output
\usepackage{lmodern}
\usepackage[T1]{fontenc}

% For cross-file-references
\usepackage{xr-hyper}

% Package for hypertext links:
\usepackage{hyperref}

% For any local file, say "hello.tex" you want to link to please
% use \externaldocument[hello-]{hello}
\externaldocument[introduction-]{introduction}
\externaldocument[conventions-]{conventions}
\externaldocument[sets-]{sets}
\externaldocument[categories-]{categories}
\externaldocument[topology-]{topology}
\externaldocument[sheaves-]{sheaves}
\externaldocument[sites-]{sites}
\externaldocument[stacks-]{stacks}
\externaldocument[fields-]{fields}
\externaldocument[algebra-]{algebra}
\externaldocument[brauer-]{brauer}
\externaldocument[homology-]{homology}
\externaldocument[derived-]{derived}
\externaldocument[simplicial-]{simplicial}
\externaldocument[more-algebra-]{more-algebra}
\externaldocument[smoothing-]{smoothing}
\externaldocument[modules-]{modules}
\externaldocument[sites-modules-]{sites-modules}
\externaldocument[injectives-]{injectives}
\externaldocument[cohomology-]{cohomology}
\externaldocument[sites-cohomology-]{sites-cohomology}
\externaldocument[dga-]{dga}
\externaldocument[dpa-]{dpa}
\externaldocument[sdga-]{sdga}
\externaldocument[hypercovering-]{hypercovering}
\externaldocument[schemes-]{schemes}
\externaldocument[constructions-]{constructions}
\externaldocument[properties-]{properties}
\externaldocument[morphisms-]{morphisms}
\externaldocument[coherent-]{coherent}
\externaldocument[divisors-]{divisors}
\externaldocument[limits-]{limits}
\externaldocument[varieties-]{varieties}
\externaldocument[topologies-]{topologies}
\externaldocument[descent-]{descent}
\externaldocument[perfect-]{perfect}
\externaldocument[more-morphisms-]{more-morphisms}
\externaldocument[flat-]{flat}
\externaldocument[groupoids-]{groupoids}
\externaldocument[more-groupoids-]{more-groupoids}
\externaldocument[etale-]{etale}
\externaldocument[chow-]{chow}
\externaldocument[intersection-]{intersection}
\externaldocument[pic-]{pic}
\externaldocument[weil-]{weil}
\externaldocument[adequate-]{adequate}
\externaldocument[dualizing-]{dualizing}
\externaldocument[duality-]{duality}
\externaldocument[discriminant-]{discriminant}
\externaldocument[derham-]{derham}
\externaldocument[local-cohomology-]{local-cohomology}
\externaldocument[algebraization-]{algebraization}
\externaldocument[curves-]{curves}
\externaldocument[resolve-]{resolve}
\externaldocument[models-]{models}
\externaldocument[functors-]{functors}
\externaldocument[equiv-]{equiv}
\externaldocument[pione-]{pione}
\externaldocument[etale-cohomology-]{etale-cohomology}
\externaldocument[proetale-]{proetale}
\externaldocument[relative-cycles-]{relative-cycles}
\externaldocument[more-etale-]{more-etale}
\externaldocument[trace-]{trace}
\externaldocument[crystalline-]{crystalline}
\externaldocument[spaces-]{spaces}
\externaldocument[spaces-properties-]{spaces-properties}
\externaldocument[spaces-morphisms-]{spaces-morphisms}
\externaldocument[decent-spaces-]{decent-spaces}
\externaldocument[spaces-cohomology-]{spaces-cohomology}
\externaldocument[spaces-limits-]{spaces-limits}
\externaldocument[spaces-divisors-]{spaces-divisors}
\externaldocument[spaces-over-fields-]{spaces-over-fields}
\externaldocument[spaces-topologies-]{spaces-topologies}
\externaldocument[spaces-descent-]{spaces-descent}
\externaldocument[spaces-perfect-]{spaces-perfect}
\externaldocument[spaces-more-morphisms-]{spaces-more-morphisms}
\externaldocument[spaces-flat-]{spaces-flat}
\externaldocument[spaces-groupoids-]{spaces-groupoids}
\externaldocument[spaces-more-groupoids-]{spaces-more-groupoids}
\externaldocument[bootstrap-]{bootstrap}
\externaldocument[spaces-pushouts-]{spaces-pushouts}
\externaldocument[spaces-chow-]{spaces-chow}
\externaldocument[groupoids-quotients-]{groupoids-quotients}
\externaldocument[spaces-more-cohomology-]{spaces-more-cohomology}
\externaldocument[spaces-simplicial-]{spaces-simplicial}
\externaldocument[spaces-duality-]{spaces-duality}
\externaldocument[formal-spaces-]{formal-spaces}
\externaldocument[restricted-]{restricted}
\externaldocument[spaces-resolve-]{spaces-resolve}
\externaldocument[formal-defos-]{formal-defos}
\externaldocument[defos-]{defos}
\externaldocument[cotangent-]{cotangent}
\externaldocument[examples-defos-]{examples-defos}
\externaldocument[algebraic-]{algebraic}
\externaldocument[examples-stacks-]{examples-stacks}
\externaldocument[stacks-sheaves-]{stacks-sheaves}
\externaldocument[criteria-]{criteria}
\externaldocument[artin-]{artin}
\externaldocument[quot-]{quot}
\externaldocument[stacks-properties-]{stacks-properties}
\externaldocument[stacks-morphisms-]{stacks-morphisms}
\externaldocument[stacks-limits-]{stacks-limits}
\externaldocument[stacks-cohomology-]{stacks-cohomology}
\externaldocument[stacks-perfect-]{stacks-perfect}
\externaldocument[stacks-introduction-]{stacks-introduction}
\externaldocument[stacks-more-morphisms-]{stacks-more-morphisms}
\externaldocument[stacks-geometry-]{stacks-geometry}
\externaldocument[moduli-]{moduli}
\externaldocument[moduli-curves-]{moduli-curves}
\externaldocument[examples-]{examples}
\externaldocument[exercises-]{exercises}
\externaldocument[guide-]{guide}
\externaldocument[desirables-]{desirables}
\externaldocument[coding-]{coding}
\externaldocument[obsolete-]{obsolete}
\externaldocument[fdl-]{fdl}
\externaldocument[index-]{index}

% Theorem environments.
%
\theoremstyle{plain}
\newtheorem{theorem}[subsection]{Theorem}
\newtheorem{proposition}[subsection]{Proposition}
\newtheorem{lemma}[subsection]{Lemma}

\theoremstyle{definition}
\newtheorem{definition}[subsection]{Definition}
\newtheorem{example}[subsection]{Example}
\newtheorem{exercise}[subsection]{Exercise}
\newtheorem{situation}[subsection]{Situation}

\theoremstyle{remark}
\newtheorem{remark}[subsection]{Remark}
\newtheorem{remarks}[subsection]{Remarks}

\numberwithin{equation}{subsection}

% Macros
%
\def\lim{\mathop{\mathrm{lim}}\nolimits}
\def\colim{\mathop{\mathrm{colim}}\nolimits}
\def\Spec{\mathop{\mathrm{Spec}}}
\def\Hom{\mathop{\mathrm{Hom}}\nolimits}
\def\Ext{\mathop{\mathrm{Ext}}\nolimits}
\def\SheafHom{\mathop{\mathcal{H}\!\mathit{om}}\nolimits}
\def\SheafExt{\mathop{\mathcal{E}\!\mathit{xt}}\nolimits}
\def\Sch{\mathit{Sch}}
\def\Mor{\mathop{\mathrm{Mor}}\nolimits}
\def\Ob{\mathop{\mathrm{Ob}}\nolimits}
\def\Sh{\mathop{\mathit{Sh}}\nolimits}
\def\NL{\mathop{N\!L}\nolimits}
\def\CH{\mathop{\mathrm{CH}}\nolimits}
\def\proetale{{pro\text{-}\acute{e}tale}}
\def\etale{{\acute{e}tale}}
\def\QCoh{\mathit{QCoh}}
\def\Ker{\mathop{\mathrm{Ker}}}
\def\Im{\mathop{\mathrm{Im}}}
\def\Coker{\mathop{\mathrm{Coker}}}
\def\Coim{\mathop{\mathrm{Coim}}}

% Boxtimes
%
\DeclareMathSymbol{\boxtimes}{\mathbin}{AMSa}{"02}

%
% Macros for moduli stacks/spaces
%
\def\QCohstack{\mathcal{QC}\!\mathit{oh}}
\def\Cohstack{\mathcal{C}\!\mathit{oh}}
\def\Spacesstack{\mathcal{S}\!\mathit{paces}}
\def\Quotfunctor{\mathrm{Quot}}
\def\Hilbfunctor{\mathrm{Hilb}}
\def\Curvesstack{\mathcal{C}\!\mathit{urves}}
\def\Polarizedstack{\mathcal{P}\!\mathit{olarized}}
\def\Complexesstack{\mathcal{C}\!\mathit{omplexes}}
% \Pic is the operator that assigns to X its picard group, usage \Pic(X)
% \Picardstack_{X/B} denotes the Picard stack of X over B
% \Picardfunctor_{X/B} denotes the Picard functor of X over B
\def\Pic{\mathop{\mathrm{Pic}}\nolimits}
\def\Picardstack{\mathcal{P}\!\mathit{ic}}
\def\Picardfunctor{\mathrm{Pic}}
\def\Deformationcategory{\mathcal{D}\!\mathit{ef}}
\usepackage{cleveref}
\crefname{originaldefinition}{original definition}{original definitions}
\Crefname{originaldefinition}{Original definition}{Original definitions}
\makeatletter
\newlabel{more-morphisms-definition-cd-morphism}{{0GTI}{source}{}{sourcepointer-1}{}}
\newlabel{more-morphisms-definition-cd-morphism@cref}{{[originaldefinition][0][]0GTI}{[1][0][]source}{}{}{}}
\makeatother
\crefname{originalequation}{original equation}{original equations}
\Crefname{originalequation}{Original equation}{Original equations}
\makeatletter
\newlabel{equation-cycle-classes}{{0H58}{source}{}{sourcepointer-2}{}}
\newlabel{equation-cycle-classes@cref}{{[originalequation][0][]0H58}{[1][0][]source}{}{}{}}
\makeatother
\crefname{originalproposition}{original proposition}{original propositions}
\Crefname{originalproposition}{Original proposition}{Original propositions}
\makeatletter
\newlabel{morphisms-proposition-generic-flatness}{{052A}{source}{}{sourcepointer-3}{}}
\newlabel{morphisms-proposition-generic-flatness@cref}{{[originalproposition][0][]052A}{[1][0][]source}{}{}{}}
\makeatother
\crefname{originallemma}{original lemma}{original lemmas}
\Crefname{originallemma}{Original lemma}{Original lemmas}
\makeatletter
\newlabel{flat-lemma-flat-after-blowing-up}{{081R}{source}{}{sourcepointer-4}{}}
\newlabel{flat-lemma-flat-after-blowing-up@cref}{{[originallemma][0][]081R}{[1][0][]source}{}{}{}}
\makeatother
\crefname{originallemma}{original lemma}{original lemmas}
\Crefname{originallemma}{Original lemma}{Original lemmas}
\makeatletter
\newlabel{lemma-uniqueness-extension}{{0H55}{source}{}{sourcepointer-5}{}}
\newlabel{lemma-uniqueness-extension@cref}{{[originallemma][0][]0H55}{[1][0][]source}{}{}{}}
\makeatother
\title{Pulling a relative cycle into the image of the cycle-classes map by a completely decomposed proper cover}
\author{The Stacks Project Authors}
\date{stacks-project:relative-cycles.tex:local-snapshot-20261012; GNU FDL 1.2 or later}
\begin{document}
\maketitle
\setcounter{subsection}{1}
\noindent\textit{Editorial note.} All delivered mathematical text is the author's own, copied from the pinned Stacks Project chapter file relative-cycles.tex; no statement or proof step was rewritten, reordered or removed. The counted span is the author environment and its own proof environment. References inside the span are handed to the assembler as navigation stubs that carry the canonical tag of the cited tag only; whether each of them is load-bearing is recorded on the stub itself.\par\medskip
\noindent\textit{Allowed standard prerequisites (editorial).}\begin{itemize}
\item Completely decomposed morphisms (More on Morphisms, Definition more-morphisms-definition-cd-morphism, canonical tag 0GTI). cited at statement lines 1180-1181. The class of morphisms g : S' to S required in the statement.
\item The cycle-classes construction (relative-cycles chapter, Equation equation-cycle-classes, canonical tag 0H58). cited at statement line 1183 and proof lines 1220 and 1229. The image of which g\textasciicircum{}*alpha must be brought.
\item Generic flatness (Morphisms, Proposition morphisms-proposition-generic-flatness, canonical tag 052A). cited at proof lines 1213-1215. Makes the closures Z\_i flat over a nonempty open U of the integral base.
\item Flatness after blowing up (More on Flatness, Lemma flat-lemma-flat-after-blowing-up, canonical tag 081R). cited at proof lines 1216-1218. Provides the U-admissible blowing up g : S' to S on which the strict transforms Z'\_i are flat over S'.
\item Uniqueness of extensions of relative cycles agreeing at the generic point (relative-cycles chapter, Lemma lemma-uniqueness-extension, canonical tag 0H55). cited at proof lines 1220-1221. Identifies beta = sum n\_i [Z'\_i/X\_\{S'\}/S']\_r with g\textasciicircum{}*alpha.
\end{itemize}\medskip
\ifdefined\hypertarget\hypertarget{sourcepointer-1}{}\fi
\noindent\textit{Source reference.} Original definition 0GTI: sources/stacks-chow/tags.txt, source lines 20968--20968. Cited at statement lines 1180-1181: the morphism class required in the statement.\par\medskip
\ifdefined\hypertarget\hypertarget{sourcepointer-2}{}\fi
\noindent\textit{Source reference.} Original equation 0H58: sources/stacks-chow/tags.txt, source lines 21378--21378. Cited at statement line 1183 and proof lines 1220 and 1229: the image of which g\textasciicircum{}*alpha is brought.\par\medskip
\ifdefined\hypertarget\hypertarget{sourcepointer-3}{}\fi
\noindent\textit{Source reference.} Original proposition 052A: sources/stacks-chow/tags.txt, source lines 6563--6563. Cited at proof lines 1213-1215: flatness of the closures over a nonempty open.\par\medskip
\ifdefined\hypertarget\hypertarget{sourcepointer-4}{}\fi
\noindent\textit{Source reference.} Original lemma 081R: sources/stacks-chow/tags.txt, source lines 10226--10226. Cited at proof lines 1216-1218: the U-admissible blowing up g : S' to S.\par\medskip
\ifdefined\hypertarget\hypertarget{sourcepointer-5}{}\fi
\noindent\textit{Source reference.} Original lemma 0H55: sources/stacks-chow/tags.txt, source lines 21375--21375. Cited at proof lines 1220-1221: identifies beta with g\textasciicircum{}*alpha.\par\medskip
\begin{lemma}
\label{lemma-get-cycles}
Let $f : X \to S$ be a finite type morphism of schemes with $S$ Noetherian.
Let $r \geq 0$. Let $\alpha$ be a relative $r$-cycle on $X/S$. Then there is
a proper, completely decomposed
(More on Morphisms, Definition \ref{more-morphisms-definition-cd-morphism})
morphism $g : S' \to S$ such that $g^*\alpha$ is in the image of
(\ref{equation-cycle-classes}).
\end{lemma}

\begin{proof}
By Noetherian induction, we may assume the result holds for the pullback of
$\alpha$ by any closed immersion $g : S' \to S$ which is not an isomorphism.

\medskip\noindent
Let $S_1 \subset S$ be an irreducible component (viewed as an integral closed
subscheme). Let $S_2 \subset S$ be the closure of the complement of $S'$
(viewed as a reduced closed subscheme). If $S_2 \not = \emptyset$, then
the result holds for the pullback of $\alpha$ by $S_1 \to S$ and $S_2 \to S$.
If $g_1 : S'_1 \to S_1$ and $g_2 : S'_2 \to S_2$
are the corresponding completely decomposed proper morphisms,
then $S' = S'_1 \amalg S'_2 \to S$
is a completely decomposed proper morphism and
we see the result holds for $S$\footnote{Namely, any closed
subscheme of $S'_1 \times_S X$ flat and of relative dimension
$\leq r$ over $S'_1$ may be viewed as a closed subscheme of $S' \times_S X$
flat and of relative dimension $\leq r$ over $S'$.}
. Thus we may assume $S' \to S$ is bijective
and we reduce to the case described in the next paragraph.

\medskip\noindent
Assume $S$ is integral. Let $\eta \in S$ be the generic point
and let $K = \kappa(\eta)$ be the function field of $S$.
Then $\alpha_\eta$ is an $r$-cycle on $X_K$.
Write $\alpha_\eta = \sum n_i[Y_i]$.
Taking the closure of $Y_i$ we obtain integral closed subschemes
$Z_i \subset X$ whose base change to $\eta$ is $Y_i$.
By generic flatness (for example Morphisms,
Proposition \ref{morphisms-proposition-generic-flatness}),
we see that $Z_i$ is flat over a nonempty open $U$ of $S$ for each $i$.
Applying More on Flatness, Lemma \ref{flat-lemma-flat-after-blowing-up}
we can find a $U$-admissible blowing up $g : S' \to S$
such that the strict transform $Z'_i \subset X_{S'}$
of $Z_i$ is flat over $S'$. Then $\beta = \sum n_i[Z'_i/X_{S'}/S']_r$
is in the image of (\ref{equation-cycle-classes}) and $\beta = g^*\alpha$
by Lemma \ref{lemma-uniqueness-extension}.

\medskip\noindent
However, this does not finish the proof as $S' \to S$ may not be
completely decomposed. This is easily fixed: denoting $T \subset S$
the complement of $U$ (viewed as a closed subscheme), by Noetherian
induction we can find a completely decomposed proper morphism
$T' \to T$ such that $(T' \to S)^*\alpha$
is in the image of (\ref{equation-cycle-classes}). Then
$S' \amalg T' \to S$ does the job.
\end{proof}

\end{document}
